\section{پرسش اول: مقایسه‌ی کابل‌های شبکه}

سرعت و کیفیت ارتباط فقط به نوع کابل بستگی ندارد.
طول کابل، کیفیت ساخت، نویز محیط و تجهیزات دو سر ارتباط هم مؤثر هستند.
با این حال، می‌توان سه نوع کابل مطرح‌شده در آزمایش را به‌صورت کلی با هم مقایسه کرد.

\subsection{کابل هم‌محور}

کابل هم‌محور\LTRfootnote{Coaxial} یک هادی مرکزی و یک لایه‌ی محافظ دارد.
این محافظ باعث می‌شود نویز محیط نسبت به کابل \لر{UTP} اثر کمتری روی سیگنال بگذارد.
در استانداردهای قدیمی اترنت، \لر{10BASE2} با سرعت ۱۰ مگابیت‌برثانیه تا ۱۸۵ متر و \لر{10BASE5} تا ۵۰۰ متر کار می‌کردند [۱].
افت سیگنال آن از فیبر بیشتر است و با زیادشدن طول و فرکانس افزایش پیدا می‌کند.

امروزه استفاده از این کابل برای یک شبکه‌ی محلی جدید معمول نیست؛
ولی اگر از قبل زیرساخت هم‌محور وجود داشته باشد، یا کاربرد مربوط به تلویزیون، آنتن و دوربین مداربسته باشد، استفاده از آن می‌تواند به‌صرفه باشد.

\subsection{کابل زوج‌به‌هم‌تابیده}

کابل زوج‌به‌هم‌تابیده\LTRfootnote{Twisted Pair} از چند زوج سیم مسی تشکیل شده است.
تابیده‌شدن سیم‌ها به دور هم، اثر نویز و هم‌شنوی را کاهش می‌دهد.
کابل \لر{Cat5e} معمولاً برای ارتباط ۱ گیگابیتی تا فاصله‌ی ۱۰۰ متر استفاده می‌شود و \لر{Cat6A} می‌تواند در همین فاصله سرعت ۱۰ گیگابیت‌برثانیه را فراهم کند [۲].

زوج‌به‌هم‌تابیده از فیبر در برابر نویز ضعیف‌تر است و در فاصله‌های زیاد افت بیشتری دارد؛
اما ارزان، در دسترس و نصب آن ساده است.
همچنین می‌توان برق تجهیزات را با \لر{PoE} روی همین کابل منتقل کرد.
به همین دلیل برای شبکه‌ی خانه، اداره، کلاس و اتصال رایانه یا نقطه‌ی دسترسی، معمولاً بهترین انتخاب است.

\subsection{فیبر نوری}

فیبر نوری اطلاعات را با نور منتقل می‌کند؛ بنابراین نویز الکترومغناطیسی روی آن اثر ندارد.
افت آن نیز بسیار کم است.
برای نمونه، استاندارد \لر{ITU-T G.652} حداکثر افت فیبر تک‌حالته را در طول موج ۱۳۱۰ نانومتر برابر ۰٫۴ دسی‌بل بر کیلومتر اعلام می‌کند [۳].
در نتیجه، فیبر برای سرعت‌های بالا و فاصله‌های طولانی مناسب‌تر است و احتمال خطای ناشی از نویز محیطی در آن کمتر خواهد بود.

تجهیزات نوری و نصب فیبر معمولاً از کابل مسی گران‌تر هستند.
به همین دلیل، استفاده از آن برای فاصله‌ی کوتاه تا یک رایانه معمولی ضروری نیست؛
ولی برای ستون فقرات شبکه، ارتباط بین دو ساختمان، مسیرهای بیشتر از ۱۰۰ متر یا محیط‌های پرنویز، هزینه‌ی بیشتر آن توجیه دارد.

در مجموع، فیبر از نظر سرعت، برد و مقاومت در برابر نویز بهترین عملکرد را دارد.
زوج‌به‌هم‌تابیده برای بیشتر شبکه‌های محلی ارزان‌ترین و ساده‌ترین گزینه است.
کابل هم‌محور نیز بیشتر زمانی مناسب است که زیرساخت آن از قبل وجود داشته باشد یا کاربرد خاصی به آن نیاز داشته باشد.

\section{پرسش دوم: معماری \لر{TCP/IP}}

معماری \لر{TCP/IP} پایه‌ی ارتباط در اینترنت است و معمولاً چهار لایه دارد [۴].

\شروع{شمارش}

\فقره
\مهم{لایه‌ی کاربرد:}
این لایه مستقیماً به برنامه‌های شبکه سرویس می‌دهد.
پروتکل‌هایی مانند \لر{HTTP}، \لر{DNS}، \لر{SMTP} و \لر{SSH} در این لایه قرار می‌گیرند.

\فقره
\مهم{لایه‌ی انتقال:}
وظیفه‌ی ارتباط بین دو برنامه را بر عهده دارد.
\لر{TCP} داده را به‌صورت مرتب و قابل‌اعتماد تحویل می‌دهد، درحالی‌که \لر{UDP} ساده‌تر است و تحویل یا ترتیب داده‌ها را تضمین نمی‌کند.

\فقره
\مهم{لایه‌ی اینترنت:}
آدرس‌دهی و رساندن بسته‌ها بین شبکه‌های مختلف در این لایه انجام می‌شود.
پروتکل اصلی این لایه \لر{IP} است و روترها با توجه به نشانی مقصد، بسته را به مسیر بعدی می‌فرستند.

\فقره
\مهم{لایه‌ی پیوند:}
این لایه ارتباط با شبکه‌ی محلی و رسانه‌ی فیزیکی را مدیریت می‌کند.
قاب اترنت، آدرس \لر{MAC} و ارسال بیت‌ها روی کابل یا فیبر در این بخش قرار می‌گیرند.

\پایان{شمارش}

مدل \لر{OSI} هفت لایه دارد، ولی \لر{TCP/IP} همین وظایف را در چهار لایه جمع می‌کند.
سه لایه‌ی کاربرد، نمایش و نشست در \لر{OSI} تقریباً در لایه‌ی کاربرد \لر{TCP/IP} قرار می‌گیرند.
لایه‌ی انتقال در هر دو مدل مشابه است.
لایه‌ی شبکه‌ی \لر{OSI} با لایه‌ی اینترنت، و دو لایه‌ی پیوند داده و فیزیکی با لایه‌ی پیوند \لر{TCP/IP} متناظر هستند.

تفاوت مهم این است که \لر{OSI} بیشتر یک مدل مرجع برای توضیح و طراحی ارتباطات شبکه است و خودش روش پیاده‌سازی مشخصی را اجبار نمی‌کند [۵].
در مقابل، \لر{TCP/IP} مجموعه‌ای از پروتکل‌های واقعی است که اینترنت با آن کار می‌کند.
پس \لر{OSI} برای درک و دسته‌بندی وظایف شبکه دقیق‌تر است، ولی \لر{TCP/IP} به چیزی که در عمل استفاده می‌شود نزدیک‌تر است.

\section{پرسش سوم: علت استفاده از کابل \لر{Straight}}

در اترنت ۱۰ و ۱۰۰ مگابیتی قدیمی، بعضی پایه‌ها مخصوص ارسال و بعضی پایه‌ها مخصوص دریافت بودند.
به همین دلیل برای اتصال دو وسیله‌ی هم‌نوع، مانند دو رایانه یا دو سوئیچ، باید جای زوج‌های ارسال و دریافت با کابل \لر{Crossover} عوض می‌شد.
برای اتصال رایانه به سوئیچ نیز کابل \لر{Straight} کافی بود.

بیشتر کارت‌های شبکه و سوئیچ‌های امروزی قابلیتی به نام \لر{Auto-MDI/MDIX} دارند.
تراشه‌ی شبکه هنگام ایجاد ارتباط تشخیص می‌دهد که کابل مستقیم است یا متقاطع.
اگر زوج‌های ارسال و دریافت جابه‌جا باشند، خود تراشه نقش آن‌ها را به‌صورت الکترونیکی عوض می‌کند [۶].
به همین دلیل، در بیشتر تجهیزات جدید هر دو نوع کابل کار می‌کنند و کاربر معمولاً فقط از کابل \لر{Straight} استفاده می‌کند.

در اترنت گیگابیتی نیز هر چهار زوج سیم به‌صورت هم‌زمان برای ارسال و دریافت به کار می‌روند [۷].
تجهیزات جدید نحوه‌ی استفاده از زوج‌ها را هنگام برقراری ارتباط تنظیم می‌کنند؛
بنابراین دیگر مانند تجهیزات قدیمی به کابل جداگانه برای اتصال وسایل هم‌نوع نیاز نیست.

البته \لر{Auto-MDI/MDIX} خرابی کابل را جبران نمی‌کند.
رنگ‌بندی دو سر کابل باید درست باشد و هیچ رشته‌ای نباید قطع یا اشتباه پرس شده باشد.
همچنین ممکن است بعضی تجهیزات قدیمی که این قابلیت را ندارند، هنوز برای اتصال مستقیم به کابل \لر{Crossover} نیاز داشته باشند.

\section{منابع}

\شروع{شمارش}

\فقره
\href{https://www.ieee802.org/misc-docs/GlobeCom2009/IEEE_802d3_Law.pdf}{\لر{IEEE 802.3, Ethernet: A Living Standard}}.

\فقره
\href{https://www.cisco.com/c/en/us/td/docs/sanity/suman/CPwE-Physical-Infrastructure_Feb2023.html}{\لر{Cisco, Physical Infrastructure Network Design for CPwE}}.

\فقره
\href{https://www.itu.int/rec/T-REC-G.652-202408-I/en}{\لر{ITU-T Recommendation G.652}}.

\فقره
\href{https://www.rfc-editor.org/rfc/rfc1122.html}{\لر{IETF RFC 1122, Requirements for Internet Hosts}}.

\فقره
\href{https://www.iso.org/standard/20269.html}{\لر{ISO/IEC 7498-1, OSI Basic Reference Model}}.

\فقره
\href{https://www.ti.com/lit/ds/symlink/dp83826ai.pdf}{\لر{Texas Instruments, DP83826 Datasheet}}.

\فقره
\href{https://www.ieee802.org/3/minutes/july98/E2_0798.pdf}{\لر{IEEE 802.3ab, 1000BASE-T Basics}}.

\پایان{شمارش}
